\section{Thiết kế chi tiết (Detailed Design)}
\label{sec:sdd_detailed_design}

\subsection{Thiết kế quy trình nghiệp vụ}
\label{subsec:sdd_business_processes}

\subsubsection{Swimlane — Quy trình sáng tác \& xuất bản}
\label{subsubsec:swim_publishing}
Swimlane Diagram được sử dụng nhằm mô tả tổng quan quy trình sáng tác và xuất bản Manga, đồng thời thể hiện sự tương tác giữa các đối tượng tham gia trong hệ thống gồm Mangaka, Assistant, Tantou Editor và Editorial Board.\\
\indent Quy trình bắt đầu từ giai đoạn đề xuất và xét duyệt Series, tiếp tục với quá trình sáng tác, phân công công việc, kiểm duyệt nội dung và kết thúc bằng việc đánh giá kết quả phát hành để đưa ra quyết định tiếp tục hoặc ngừng xuất bản Series. 
\begin{figure}[H]
    \centering
    \includegraphics[height=0.75\textheight, width=\textwidth, keepaspectratio]{Chapter3_IMAGE/Swimlane_tong_quat_quy_trinh_sang_tac_maga.png}
    \caption{Swimlane Diagram tổng quát quy trình sáng tác và xuất bản Manga}
    \label{fig:swimlane_manga}
\end{figure}

\FloatBarrier

\subsubsection{Activity Diagram — Giao và thực hiện Task}
\label{subsubsec:act_task}
Activity Diagram này mô tả quy trình phân công và thực hiện công việc giữa Mangaka và Assistant.\\
\indent Quy trình bắt đầu khi Mangaka tạo Task và giao cho Assistant. Sau khi nhận nhiệm vụ, Assistant thực hiện công việc, gửi Submission lên hệ thống và chờ kết quả kiểm duyệt. Trong trường hợp Submission chưa đạt yêu cầu, Assistant sẽ tiếp tục chỉnh sửa và gửi lại cho đến khi được chấp thuận.
\begin{figure}[H]
    \centering
    \includegraphics[height=0.75\textheight, width=0.8\textwidth, keepaspectratio]{Chapter3_IMAGE/Activity_Diagram_quy_trinh_giao_va_thuc_hien_Task.png}
    \caption{Activity Diagram quy trình giao và thực hiện Task}
    \label{fig:activity_task}
\end{figure}

\FloatBarrier

\subsubsection{Activity Diagram — Review Chapter}
\label{subsubsec:act_review}
Activity Diagram này mô tả quy trình kiểm duyệt nội dung Manga giữa Mangaka và Tantou Editor.\\
\indent Sau khi Chapter được hoàn thiện, Tantou Editor tiến hành xem xét nội dung và đưa ra phản hồi. Nếu Chapter chưa đạt yêu cầu, Mangaka sẽ thực hiện chỉnh sửa theo nhận xét. Quy trình được lặp lại cho đến khi Chapter được phê duyệt.
\begin{figure}[H]
    \centering
    \includegraphics[height=0.75\textheight, width=0.8\textwidth, keepaspectratio]{Chapter3_IMAGE/Activity_Diagram_quy_trinh_Review_Chapter.png}
    \caption{Activity Diagram quy trình Review Chapter}
    \label{fig:activity_review}
\end{figure}

\FloatBarrier

\subsubsection{Activity Diagram — Xuất bản Manga}
\label{subsubsec:act_publish}
Activity Diagram này mô tả quá trình đánh giá và ra quyết định xuất bản đối với một Series Manga.\\
\indent Editorial Board thực hiện cập nhật dữ liệu bình chọn, theo dõi thứ hạng Series và đánh giá hiệu quả phát hành. Dựa trên kết quả đánh giá, hội đồng đưa ra quyết định tiếp tục xuất bản, điều chỉnh kế hoạch phát hành hoặc ngừng phát hành Series.
\begin{figure}[H]
    \centering
    \includegraphics[height=0.75\textheight, width=0.8\textwidth, keepaspectratio]{Chapter3_IMAGE/Activity_Diagram_quy_trinh_xuat_ban_Manga.png}
    \caption{Activity Diagram quy trình xuất bản Manga}
    \label{fig:activity_xuat_ban}
\end{figure}

\FloatBarrier

\subsection{Thiết kế Sequence Diagram}
\label{subsec:sdd_sequence_diagrams}

\subsubsection{Sequence — Đăng nhập hệ thống}
\label{subsubsec:seq_login}
Sequence Diagram mô tả quá trình người dùng đăng nhập vào hệ thống. Hệ thống xác thực thông tin tài khoản và cấp quyền truy cập nếu hợp lệ.

\begin{figure}[H]
    \centering
    \includegraphics[width=\textwidth]{Chapter3_IMAGE/sequence_diagram_dang_nhap_he_thong.png}
    \caption{Sequence Diagram Đăng nhập hệ thống}
    \label{fig:seq_dang_nhap}
\end{figure}

\subsubsection{Sequence — Phân công Task}
\label{subsubsec:seq_task}
Sequence Diagram mô tả quá trình Manager tạo và phân công công việc cho Assistant. Sau khi Task được tạo, hệ thống lưu thông tin và gửi thông báo đến Assistant.

\begin{figure}[H]
    \centering
    \includegraphics[width=\textwidth]{Chapter3_IMAGE/sequence_diagram_manga_task_assignment.png}
    \caption{Sequence Diagram Phân công Task}
    \label{fig:seq_giao_task}
\end{figure}

\subsubsection{Sequence — Nộp Submission}
\label{subsubsec:seq_submission}
Sequence Diagram mô tả quá trình Assistant nộp kết quả thực hiện Task. Hệ thống lưu Submission và cập nhật trạng thái công việc.

\begin{figure}[H]
    \centering
    \includegraphics[width=\textwidth]{Chapter3_IMAGE/sequence_diagram_quy_trinh_nop_submission.png}
    \caption{Sequence Diagram Nộp Submission}
    \label{fig:seq_nop_submission}
\end{figure}

\subsubsection{Sequence — Review Chapter}
\label{subsubsec:seq_review}
Sequence Diagram mô tả quá trình Manager đánh giá Submission. Nếu đạt yêu cầu, hệ thống cập nhật tiến độ Chapter; ngược lại hệ thống gửi yêu cầu chỉnh sửa.

\begin{figure}[H]
    \centering
    \includegraphics[width=\textwidth]{Chapter3_IMAGE/sequence_diagram_review_chapter.png}
    \caption{Sequence Diagram Review Chapter}
    \label{fig:seq_review_chapter}
\end{figure}

\subsubsection{Sequence — Phê duyệt Xuất bản Series}
\label{subsubsec:seq_publish}
Sequence Diagram mô tả quá trình Hội đồng Biên tập (Editorial Board) tiến hành bỏ phiếu và phê duyệt bộ truyện mới (Series). Hệ thống tự động theo dõi tỷ lệ tán thành và chỉ chuyển trạng thái bộ truyện sang Đang triển khai (Ongoing) khi 100% thành viên Hội đồng đã bỏ phiếu và tỷ lệ tán thành đạt từ 50% trở lên. Sau khi phê duyệt, hệ thống sẽ ghi nhận nhật ký, cập nhật lịch xuất bản, phân công Tantou Editor và gửi thông báo kết quả cho Mangaka.

\begin{figure}[H]
    \centering
    \includegraphics[width=\textwidth]{Chapter3_IMAGE/sequence_diagram_manga_series_publishing.png}
    \caption{Sequence Diagram Phê duyệt Xuất bản Series}
    \label{fig:seq_xuat_ban}
\end{figure}

\subsection{Thiết kế các lớp học}
\label{subsec:sdd_class_diagram}

Class Diagram mô tả cấu trúc tĩnh của hệ thống, bao gồm các lớp chính, thuộc tính, phương thức và mối quan hệ giữa các lớp trong hệ thống quản lý quy trình sản xuất Manga.

\subsubsection{Chi tiết các lớp trong hệ thống}
\label{subsubsec:class_details}

\subsubsection{Lớp Role}
\textit{Mô tả:} Quản lý thông tin vai trò của người dùng trong hệ thống.

\begin{table}[H]
\centering
\renewcommand{\arraystretch}{1.3}
\begin{tabularx}{\textwidth}{|p{4cm}|p{3cm}|>{\raggedright\arraybackslash}X|}
\hline
\rowcolor{tableheader}
\textbf{Thuộc tính} & \textbf{Kiểu dữ liệu} & \textbf{Ý nghĩa} \\ \hline
role\_id & int & Mã vai trò \\ \hline
role\_name & string & Tên vai trò \\ \hline
description & string & Mô tả vai trò \\ \hline
\end{tabularx}
\caption{Đặc tả chi tiết thuộc tính lớp Role}
\label{tab:class_role}
\end{table}


\subsubsection{Lớp User}
\textit{Mô tả:} Quản lý thông tin tài khoản người dùng.

\begin{table}[H]
\centering
\renewcommand{\arraystretch}{1.3}
\begin{tabularx}{\textwidth}{|p{4cm}|p{3cm}|>{\raggedright\arraybackslash}X|}
\hline
\rowcolor{tableheader}
\textbf{Thuộc tính} & \textbf{Kiểu dữ liệu} & \textbf{Ý nghĩa} \\ \hline
user\_id & int & Mã người dùng (Primary Key) \\ \hline
username & string & Tên đăng nhập \\ \hline
full\_name & string & Họ và tên \\ \hline
email & string & Địa chỉ email \\ \hline
password\_hash & string & Mật khẩu đã mã hóa \\ \hline
status & string & Trạng thái tài khoản \\ \hline
is\_head\_board & int & Đánh dấu thành viên là trưởng ban hội đồng xét duyệt \\ \hline
role\_id & int & Mã vai trò (Foreign Key liên kết lớp Role) \\ \hline
\end{tabularx}
\caption{Đặc tả chi tiết thuộc tính lớp User}
\label{tab:class_user}
\end{table}


\subsubsection{Lớp Series}
\textit{Mô tả:} Quản lý thông tin bộ truyện Manga.

\begin{table}[H]
\centering
\renewcommand{\arraystretch}{1.3}
\begin{tabularx}{\textwidth}{|p{4cm}|p{3cm}|>{\raggedright\arraybackslash}X|}
\hline
\rowcolor{tableheader}
\textbf{Thuộc tính} & \textbf{Kiểu dữ liệu} & \textbf{Ý nghĩa} \\ \hline
series\_id & int & Mã bộ truyện (Primary Key) \\ \hline
title & string & Tên bộ truyện \\ \hline
description & string & Mô tả nội dung \\ \hline
status & string & Trạng thái bộ truyện (\texttt{'planning'}, \texttt{'ongoing'}...) \\ \hline
publish\_type & string & Tần suất phát hành (weekly, monthly...) \\ \hline
cover\_image & string & Ảnh bìa bộ truyện \\ \hline
proposal\_file & string & Đường dẫn tệp đề xuất bộ truyện \\ \hline
editor\_id & int & Mã Biên tập viên phụ trách (Foreign Key, cho phép NULL) \\ \hline
dossier\_notes & string & Ghi chú hồ sơ duyệt bộ truyện \\ \hline
mangaka\_id & int & Mã tác giả sở hữu (Foreign Key liên kết lớp User) \\ \hline
\end{tabularx}
\caption{Đặc tả chi tiết thuộc tính lớp Series}
\label{tab:class_series}
\end{table}


\subsubsection{Lớp Chapter}
\textit{Mô tả:} Quản lý thông tin chương truyện.

\begin{table}[H]
\centering
\renewcommand{\arraystretch}{1.3}
\begin{tabularx}{\textwidth}{|p{4cm}|p{3cm}|>{\raggedright\arraybackslash}X|}
\hline
\rowcolor{tableheader}
\textbf{Thuộc tính} & \textbf{Kiểu dữ liệu} & \textbf{Ý nghĩa} \\ \hline
chapter\_id & int & Mã chương (Primary Key) \\ \hline
chapter\_number & int & Số thứ tự chương \\ \hline
title & string & Tên chương \\ \hline
status & string & Trạng thái chương \\ \hline
due\_date & datetime & Hạn hoàn thành chương truyện \\ \hline
published\_at & datetime & Ngày xuất bản \\ \hline
is\_final & tinyint & Đánh dấu chương cuối cùng của bộ truyện \\ \hline
series\_id & int & Mã bộ truyện (Foreign Key liên kết lớp Series) \\ \hline
\end{tabularx}
\caption{Đặc tả chi tiết thuộc tính lớp Chapter}
\label{tab:class_chapter}
\end{table}


\subsubsection{Lớp Page}
\textit{Mô tả:} Quản lý các trang truyện thuộc Chapter.

\begin{table}[H]
\centering
\renewcommand{\arraystretch}{1.3}
\begin{tabularx}{\textwidth}{|p{4cm}|p{3cm}|>{\raggedright\arraybackslash}X|}
\hline
\rowcolor{tableheader}
\textbf{Thuộc tính} & \textbf{Kiểu dữ liệu} & \textbf{Ý nghĩa} \\ \hline
page\_id & int & Mã trang (Primary Key) \\ \hline
page\_number & int & Số thứ tự trang \\ \hline
image\_url & string & Đường dẫn ảnh \\ \hline
old\_image\_url & string & Đường dẫn ảnh cũ trước khi sửa đổi \\ \hline
status & string & Trạng thái hoàn thành \\ \hline
chapter\_id & int & Mã chương (Foreign Key liên kết lớp Chapter) \\ \hline
\end{tabularx}
\caption{Đặc tả chi tiết thuộc tính lớp Page}
\label{tab:class_page}
\end{table}


\subsubsection{Lớp PageRegion}
\textit{Mô tả:} Lưu thông tin phân vùng cụ thể của trang vẽ để phân công công việc.

\begin{table}[H]
\centering
\renewcommand{\arraystretch}{1.3}
\begin{tabularx}{\textwidth}{|p{4cm}|p{3cm}|>{\raggedright\arraybackslash}X|}
\hline
\rowcolor{tableheader}
\textbf{Thuộc tính} & \textbf{Kiểu dữ liệu} & \textbf{Ý nghĩa} \\ \hline
region\_id & int & Mã phân vùng (Primary Key) \\ \hline
region\_type & string & Loại phân vùng (panel, bubble, background...) \\ \hline
x & int & Tọa độ x trên trang truyện \\ \hline
y & int & Tọa độ y trên trang truyện \\ \hline
width & int & Chiều rộng phân vùng \\ \hline
height & int & Chiều cao phân vùng \\ \hline
status & string & Trạng thái hoàn thành phân vùng \\ \hline
page\_id & int & Mã trang truyện chứa phân vùng (Foreign Key liên kết lớp Page) \\ \hline
\end{tabularx}
\caption{Đặc tả chi tiết thuộc tính lớp PageRegion}
\label{tab:class_pageregion}
\end{table}


\subsubsection{Lớp Task}
\textit{Mô tả:} Quản lý công việc được giao cho Assistant.

\begin{table}[H]
\centering
\renewcommand{\arraystretch}{1.3}
\begin{tabularx}{\textwidth}{|p{4cm}|p{3cm}|>{\raggedright\arraybackslash}X|}
\hline
\rowcolor{tableheader}
\textbf{Thuộc tính} & \textbf{Kiểu dữ liệu} & \textbf{Ý nghĩa} \\ \hline
task\_id & int & Mã công việc (Primary Key) \\ \hline
title & string & Tên công việc \\ \hline
task\_type & string & Loại công việc (background, inking, coloring...) \\ \hline
description & string & Mô tả công việc \\ \hline
resource\_url & string & Đường dẫn tệp tài nguyên đính kèm \\ \hline
priority & string & Mức độ ưu tiên \\ \hline
status & string & Trạng thái công việc \\ \hline
due\_date & datetime & Hạn hoàn thành \\ \hline
page\_id & int & Mã trang truyện liên kết (Foreign Key liên kết lớp Page) \\ \hline
page\_region\_id & int & Mã phân vùng liên quan (Foreign Key, cho phép NULL) \\ \hline
grouped\_region\_ids & string & Danh sách các phân vùng gộp liên quan (cho phép NULL) \\ \hline
mangaka\_id & int & Mã tác giả giao việc (Foreign Key liên kết lớp User) \\ \hline
assistant\_id & int & Mã trợ lý nhận việc (Foreign Key liên kết lớp User) \\ \hline
\end{tabularx}
\caption{Đặc tả chi tiết thuộc tính lớp Task}
\label{tab:class_task}
\end{table}


\subsubsection{Lớp Submission}
\textit{Mô tả:} Quản lý các sản phẩm được nộp lên hệ thống.

\begin{table}[H]
\centering
\renewcommand{\arraystretch}{1.3}
\begin{tabularx}{\textwidth}{|p{4cm}|p{3cm}|>{\raggedright\arraybackslash}X|}
\hline
\rowcolor{tableheader}
\textbf{Thuộc tính} & \textbf{Kiểu dữ liệu} & \textbf{Ý nghĩa} \\ \hline
submission\_id & int & Mã bài nộp (Primary Key) \\ \hline
file\_url & string & Đường dẫn tập tin \\ \hline
notes & string & Ghi chú bổ sung \\ \hline
status & string & Trạng thái bài nộp \\ \hline
submitted\_at & datetime & Thời gian nộp \\ \hline
task\_id & int & Mã công việc liên kết (Foreign Key liên kết lớp Task, cho phép NULL) \\ \hline
chapter\_id & int & Mã chương liên kết (Foreign Key liên kết lớp Chapter, cho phép NULL) \\ \hline
user\_id & int & Mã người nộp (Foreign Key liên kết lớp User) \\ \hline
\end{tabularx}
\caption{Đặc tả chi tiết thuộc tính lớp Submission}
\label{tab:class_submission}
\end{table}


\subsubsection{Lớp Review}
\textit{Mô tả:} Quản lý các đánh giá và phản hồi đối với Submission.

\begin{table}[H]
\centering
\renewcommand{\arraystretch}{1.3}
\begin{tabularx}{\textwidth}{|p{4cm}|p{3cm}|>{\raggedright\arraybackslash}X|}
\hline
\rowcolor{tableheader}
\textbf{Thuộc tính} & \textbf{Kiểu dữ liệu} & \textbf{Ý nghĩa} \\ \hline
review\_id & int & Mã đánh giá (Primary Key) \\ \hline
comments & string & Nội dung nhận xét \\ \hline
rating & int & Điểm đánh giá \\ \hline
created\_at & datetime & Thời gian đánh giá \\ \hline
submission\_id & int & Mã bài nộp được đánh giá (Foreign Key liên kết lớp Submission) \\ \hline
reviewer\_id & int & Mã người đánh giá (Foreign Key liên kết lớp User) \\ \hline
\end{tabularx}
\caption{Đặc tả chi tiết thuộc tính lớp Review}
\label{tab:class_review}
\end{table}


\subsubsection{Lớp EditorAnnotation}
\textit{Mô tả:} Lưu thông tin ghi chú lỗi trực quan vẽ đè trên trang truyện từ Editor.

\begin{table}[H]
\centering
\renewcommand{\arraystretch}{1.3}
\begin{tabularx}{\textwidth}{|p{4cm}|p{3cm}|>{\raggedright\arraybackslash}X|}
\hline
\rowcolor{tableheader}
\textbf{Thuộc tính} & \textbf{Kiểu dữ liệu} & \textbf{Ý nghĩa} \\ \hline
annotation\_id & int & Mã ghi chú (Primary Key) \\ \hline
x & int & Tọa độ x của vùng đánh dấu nét lỗi \\ \hline
y & int & Tọa độ y của vùng đánh dấu nét lỗi \\ \hline
width & int & Chiều rộng vùng lỗi \\ \hline
height & int & Chiều cao vùng lỗi \\ \hline
comments & string & Nhận xét, góp ý chi tiết \\ \hline
created\_at & datetime & Thời gian tạo ghi chú \\ \hline
page\_id & int & Mã trang truyện liên quan (Foreign Key liên kết lớp Page) \\ \hline
editor\_id & int & Mã Biên tập viên ghi chú (Foreign Key liên kết lớp User) \\ \hline
\end{tabularx}
\caption{Đặc tả chi tiết thuộc tính lớp EditorAnnotation}
\label{tab:class_editorannotation}
\end{table}


\subsubsection{Lớp SubmissionAnnotation}
\textit{Mô tả:} Lưu trữ thông tin ghi chú báo lỗi trực quan vẽ đè của Tác giả (Mangaka) trên các bản vẽ nộp của Trợ lý.

\begin{table}[H]
\centering
\renewcommand{\arraystretch}{1.3}
\begin{tabularx}{\textwidth}{|p{4cm}|p{3cm}|>{\raggedright\arraybackslash}X|}
\hline
\rowcolor{tableheader}
\textbf{Thuộc tính} & \textbf{Kiểu dữ liệu} & \textbf{Ý nghĩa} \\ \hline
annotation\_id & int & Mã ghi chú lỗi (Primary Key) \\ \hline
x & int & Tọa độ x của vùng báo lỗi trên hình ảnh nộp \\ \hline
y & int & Tọa độ y của vùng báo lỗi trên hình ảnh nộp \\ \hline
width & int & Chiều rộng vùng báo lỗi \\ \hline
height & int & Chiều cao vùng báo lỗi \\ \hline
comments & string & Ghi chú mô tả lỗi chi tiết \\ \hline
created\_at & datetime & Thời gian tạo ghi chú \\ \hline
submission\_id & int & Mã bản nộp của Trợ lý liên quan (Foreign Key liên kết lớp Submission) \\ \hline
user\_id & int & Mã Tác giả thực hiện ghi chú (Foreign Key liên kết lớp User) \\ \hline
\end{tabularx}
\caption{Đặc tả chi tiết thuộc tính lớp SubmissionAnnotation}
\label{tab:class_submissionannotation}
\end{table}


\subsubsection{Lớp BoardVote}
\textit{Mô tả:} Lưu trữ phiếu biểu quyết của thành viên Hội đồng biên tập cho đề xuất Series.

\begin{table}[H]
\centering
\renewcommand{\arraystretch}{1.3}
\begin{tabularx}{\textwidth}{|p{4cm}|p{3cm}|>{\raggedright\arraybackslash}X|}
\hline
\rowcolor{tableheader}
\textbf{Thuộc tính} & \textbf{Kiểu dữ liệu} & \textbf{Ý nghĩa} \\ \hline
vote\_id & int & Mã phiếu bầu (Primary Key) \\ \hline
vote & string & Quyết định bầu chọn (approve/reject) \\ \hline
created\_at & datetime & Thời gian bỏ phiếu \\ \hline
series\_id & int & Mã bộ truyện liên quan (Foreign Key liên kết lớp Series) \\ \hline
board\_member\_id & int & Mã thành viên hội đồng bỏ phiếu (Foreign Key liên kết lớp User) \\ \hline
\end{tabularx}
\caption{Đặc tả chi tiết thuộc tính lớp BoardVote}
\label{tab:class_boardvote}
\end{table}


\subsubsection{Lớp SeriesRanking}
\textit{Mô tả:} Quản lý thông tin xếp hạng của các bộ truyện.

\begin{table}[H]
\centering
\renewcommand{\arraystretch}{1.3}
\begin{tabularx}{\textwidth}{|p{4cm}|p{3cm}|>{\raggedright\arraybackslash}X|}
\hline
\rowcolor{tableheader}
\textbf{Thuộc tính} & \textbf{Kiểu dữ liệu} & \textbf{Ý nghĩa} \\ \hline
ranking\_id & int & Mã xếp hạng (Primary Key) \\ \hline
rank\_position & int & Thứ hạng \\ \hline
score & decimal & Điểm đánh giá \\ \hline
votes & int & Số lượng phiếu bầu trong kỳ \\ \hline
period\_start\_date & date & Thời điểm bắt đầu kỳ xếp hạng \\ \hline
series\_id & int & Mã bộ truyện được xếp hạng (Foreign Key liên kết lớp Series) \\ \hline
board\_member\_id & int & Mã thành viên BBT chấm điểm (Foreign Key liên kết lớp User) \\ \hline
\end{tabularx}
\caption{Đặc tả chi tiết thuộc tính lớp SeriesRanking}
\label{tab:class_seriesranking}
\end{table}


\subsubsection{Lớp Notification}
\textit{Mô tả:} Quản lý thông báo gửi đến người dùng.

\begin{table}[H]
\centering
\renewcommand{\arraystretch}{1.3}
\begin{tabularx}{\textwidth}{|p{4cm}|p{3cm}|>{\raggedright\arraybackslash}X|}
\hline
\rowcolor{tableheader}
\textbf{Thuộc tính} & \textbf{Kiểu dữ liệu} & \textbf{Ý nghĩa} \\ \hline
notification\_id & int & Mã thông báo (Primary Key) \\ \hline
type & string & Loại thông báo \\ \hline
message & string & Nội dung thông báo \\ \hline
is\_read & boolean & Trạng thái đã đọc \\ \hline
related\_id & int & ID của thực thể liên quan (task, submission...) \\ \hline
user\_id & int & Mã người dùng nhận (Foreign Key liên kết lớp User) \\ \hline
\end{tabularx}
\caption{Đặc tả chi tiết thuộc tính lớp Notification}
\label{tab:class_notification}
\end{table}


\subsubsection{Lớp SystemLog}
\textit{Mô tả:} Lưu thông tin nhật ký các hoạt động nhạy cảm trong hệ thống.

\begin{table}[H]
\centering
\renewcommand{\arraystretch}{1.3}
\begin{tabularx}{\textwidth}{|p{4cm}|p{3cm}|>{\raggedright\arraybackslash}X|}
\hline
\rowcolor{tableheader}
\textbf{Thuộc tính} & \textbf{Kiểu dữ liệu} & \textbf{Ý nghĩa} \\ \hline
log\_id & int & Mã log (Primary Key) \\ \hline
action & string & Tên hành động thực hiện \\ \hline
details & string & Mô tả chi tiết \\ \hline
ip\_address & string & Địa chỉ IP truy cập \\ \hline
created\_at & datetime & Thời gian ghi nhận log \\ \hline
user\_id & int & Mã người dùng thực hiện (Foreign Key liên kết lớp User, có thể NULL) \\ \hline
\end{tabularx}
\caption{Đặc tả chi tiết thuộc tính lớp SystemLog}
\label{tab:class_systemlog}
\end{table}


\subsubsection{Bổ sung phương thức (Methods)}
\label{subsubsec:class_methods}

Các phương thức (Methods) của từng lớp trong hệ thống mô tả chi tiết chữ ký hàm (Method Signatures) bao gồm tham số truyền vào, kiểu dữ liệu trả về và ý nghĩa hoạt động thực tế trong mã nguồn PHP Model:

\begin{itemize}
    \item \textbf{Lớp Role:}
    \begin{itemize}
        \item \texttt{findByRoleName(roleName: string): Role}\\
            Lấy thông tin vai trò theo tên (ví dụ: 'mangaka', 'assistant'...).
        \item \texttt{getRolesWithUserCount(): array}\\
            Lấy danh sách vai trò kèm theo số lượng người dùng thuộc vai trò đó.
    \end{itemize}
    
    \item \textbf{Lớp User:}
    \begin{itemize}
        \item \texttt{findByEmail(email: string): array}\\
            Tìm người dùng theo email đăng ký.
        \item \texttt{findByUsername(username: string): array}\\
            Tìm người dùng theo tên đăng nhập.
        \item \texttt{getUserByIdWithRole(id: int): array}\\
            Lấy thông tin người dùng kèm theo tên vai trò tương ứng.
        \item \texttt{findByRoleName(roleName: string): array}\\
            Lấy danh sách người dùng đang active theo tên vai trò.
        \item \texttt{getRoleStatistics(): array}\\
            Lấy dữ liệu thống kê số lượng tài khoản theo từng vai trò.
        \item \texttt{getAllUsersWithRole(): array}\\
            Lấy danh sách tất cả người dùng trong hệ thống cùng vai trò của họ.
        \item \texttt{getPaginatedUsers(limit: int, offset: int): array}\\
            Lấy danh sách người dùng phân trang phục vụ trang quản trị Admin.
        \item \texttt{demoteOtherHeads(exceptUserId: int): bool}\\
            Bãi nhiệm vai trò Trưởng ban hội đồng của các thành viên khác khi có Trưởng ban mới được bổ nhiệm.
    \end{itemize}

    \item \textbf{Lớp Series:}
    \begin{itemize}
        \item \texttt{findByMangakaId(mangakaId: int): array}\\
            Lấy danh sách các bộ truyện của tác giả đang đăng nhập.
        \item \texttt{findByStatus(status: string): array}\\
            Lấy danh sách các bộ truyện theo trạng thái (ongoing, planning...).
        \item \texttt{getSeriesByEditorId(editorId: int): array}\\
            Lấy danh sách các bộ truyện được phân công cho Biên tập viên và đã duyệt phát hành.
        \item \texttt{getSeriesForPublishing(): array}\\
            Lấy danh sách bộ truyện kèm thống kê số chương, điểm số phục vụ cho Hội đồng duyệt xuất bản.
        \item \texttt{getDossiersByEditorId(editorId: int): array}\\
            Lấy danh sách hồ sơ truyện theo Biên tập viên phụ trách.
        \item \texttt{countActiveSeries(): int}\\
            Đếm số lượng bộ truyện đang hoạt động trong hệ thống.
        \item \texttt{isTitleExists(title: string, excludeId: int = null): bool}\\
            Kiểm tra xem tên truyện đã tồn tại trên hệ thống chưa để tránh trùng tên.
    \end{itemize}

    \item \textbf{Lớp Chapter:}
    \begin{itemize}
        \item \texttt{findBySeriesId(seriesId: int): array}\\
            Lấy danh sách các chương thuộc bộ truyện, sắp xếp tăng dần.
        \item \texttt{isChapterNumberExists(seriesId: int, chapterNumber: int): bool}\\
            Kiểm tra số chương đã tồn tại trong bộ truyện chưa để tránh trùng lặp.
        \item \texttt{findByMangakaId(mangakaId: int): array}\\
            Lấy danh sách chương truyện của Mangaka sở hữu.
        \item \texttt{findBySeriesIdWithStats(seriesId: int): array}\\
            Lấy danh sách chương truyện kèm thống kê số lượng Task hoàn thành của trợ lý.
        \item \texttt{hasFinalChapter(seriesId: int): bool}\\
            Kiểm tra xem truyện đã có chương cuối cùng (is\_final = 1) chưa.
        \item \texttt{countUnfinishedChapters(seriesId: int): int}\\
            Đếm số chương chưa hoàn tất của một truyện.
        \item \texttt{hasFinalApprovedChapter(seriesId: int): bool}\\
            Kiểm tra xem chương cuối đã được duyệt hay chưa.
        \item \texttt{getUpcomingDeadlinesByEditor(editorId: int): array}\\
            Lấy danh sách các chương sắp đến hạn nộp thuộc BTV quản lý.
        \item \texttt{findPublishedChapters(): array}\\
            Lấy danh sách các chương đã xuất bản (published) thuộc các Series.
        \item \texttt{findByMangakaIdWithStats(mangakaId: int): array}\\
            Lấy danh sách chương truyện kèm thống kê số lượng Task hoàn thành và nộp bài cho Mangaka.
        \item \texttt{countByMangakaId(mangakaId: int): int}\\
            Đếm tổng số chương truyện do Mangaka quản lý/sở hữu.
        \item \texttt{findApprovedChapters(): array}\\
            Lấy tất cả các chương truyện đã ở trạng thái được phê duyệt hoàn thiện (approved).
    \end{itemize}

    \item \textbf{Lớp Page:}
    \begin{itemize}
        \item \texttt{findByChapterId(chapterId: int): array}\\
            Lấy danh sách các trang vẽ thuộc chương, sắp xếp theo số trang.
        \item \texttt{isPageNumberExists(chapterId: int, pageNumber: int): bool}\\
            Kiểm tra trùng số trang trong một chương.
        \item \texttt{findByChapterIdWithAnnotationCount(chapterId: int): array}\\
            Lấy danh sách trang vẽ kèm số lỗi ghi chú của Biên tập viên.
        \item \texttt{getPageNumbersByChapterId(chapterId: int): array}\\
            Lấy mảng các số thứ tự trang hiện có trong chương.
        \item \texttt{getOtherPageNumbers(chapterId: int, pageId: int): array}\\
            Lấy số thứ tự các trang khác trong chương để phục vụ thay đổi vị trí.
        \item \texttt{updateStatusByChapterId(chapterId: int, status: string): bool}\\
            Cập nhật trạng thái cho tất cả các trang thuộc một chương truyện cụ thể.
        \item \texttt{findByMangakaId(mangakaId: int): array}\\
            Lấy danh sách các trang vẽ do trợ lý nộp thuộc quyền quản lý của Mangaka.
        \item \texttt{countByMangakaId(mangakaId: int): int}\\
            Đếm số lượng trang vẽ do Mangaka tạo ra trên hệ thống.
    \end{itemize}

    \item \textbf{Lớp PageRegion:}
    \begin{itemize}
        \item \texttt{findByPageId(pageId: int): array}\\
            Lấy tất cả các phân vùng của một trang truyện cụ thể.
        \item \texttt{findByIds(ids: array): array}\\
            Tìm các phân vùng trên trang dựa theo danh sách ID truyền vào.
    \end{itemize}

    \item \textbf{Lớp Task:}
    \begin{itemize}
        \item \texttt{findByPageId(pageId: int): array}\\
            Lấy danh sách các công việc trên một trang truyện cụ thể kèm tên trợ lý.
        \item \texttt{findByAssistantId(assistantId: int): array}\\
            Lấy danh sách công việc được giao của trợ lý kèm thông tin chương, trang, phân vùng.
        \item \texttt{findActiveByAssistantId(assistantId: int): array}\\
            Lấy các công việc chưa hoàn thành của một trợ lý.
        \item \texttt{findByMangakaId(mangakaId: int): array}\\
            Lấy danh sách công việc do tác giả giao cho các trợ lý.
        \item \texttt{getStatusStatistics(): array}\\
            Lấy thống kê số lượng Task theo trạng thái (pending, in\_progress, completed).
        \item \texttt{getMonthlyIncomeStats(assistantId: int): array}\\
            Tính toán thống kê thu nhập dự kiến của trợ lý theo từng tháng dựa trên số Task hoàn thành.
        \item \texttt{findTasksByChapterId(chapterId: int): array}\\
            Tìm danh sách toàn bộ Task liên kết với các trang trong một Chapter.
        \item \texttt{findByPageRegionIds(regionIds: array): array}\\
            Tìm các Task liên quan đến danh sách ID phân vùng.
        \item \texttt{deleteActiveTasksByRegionIds(regionIds: array): bool}\\
            Xóa nhanh các Task đang hoạt động liên kết với danh sách phân vùng.
        \item \texttt{countByPageAndAssistant(pageId: int, assistantId: int): int}\\
            Đếm số Task được giao cho trợ lý cụ thể trên một trang truyện.
        \item \texttt{countByAssistantIdWithGating(assistantId: int): int}\\
            Đếm số Task trợ lý đang thực hiện để thực hiện giới hạn nhận việc tối đa (Gating).
    \end{itemize}

    \item \textbf{Lớp Submission:}
    \begin{itemize}
        \item \texttt{findByTaskId(taskId: int): array}\\
            Lấy lịch sử các lần nộp bài vẽ cho một Task cụ thể.
        \item \texttt{findByChapterId(chapterId: int): array}\\
            Lấy danh sách các lần nộp bản thảo cho một chương cụ thể.
        \item \texttt{findPendingSubmissions(): array}\\
            Lấy danh sách các bản thảo chương đang chờ duyệt của hệ thống.
        \item \texttt{countPendingChapterSubmissionsByEditor(editorId: int): int}\\
            Đếm số lượng chương đang chờ duyệt thuộc Biên tập viên quản lý.
        \item \texttt{getStatusStatistics(): array}\\
            Lấy thống kê số lượng bài nộp theo trạng thái.
        \item \texttt{countByMangakaId(mangakaId: int): int}\\
            Đếm tổng số bài nộp do Trợ lý thực hiện dưới quyền của Mangaka.
        \item \texttt{countPendingByMangakaId(mangakaId: int): int}\\
            Đếm số bài nộp đang chờ duyệt thuộc quyền quản lý của Mangaka.
        \item \texttt{findByUserId(userId: int): array}\\
            Lấy danh sách các bản nộp dựa theo ID người nộp (User ID).
        \item \texttt{findPendingSubmissionsByEditorId(editorId: int): array}\\
            Lấy các bản nộp chương/task đang chờ Biên tập viên chuyên trách phê duyệt.
        \item \texttt{findAllChapterSubmissionsByEditorId(editorId: int): array}\\
            Lấy toàn bộ lịch sử bản nộp chương liên kết với Biên tập viên.
        \item \texttt{findAllTaskSubmissionsByMangakaId(mangakaId: int): array}\\
            Lấy toàn bộ lịch sử nộp bài vẽ phân vùng của Trợ lý cho Mangaka.
        \item \texttt{findWithMetadataById(id: int): array}\\
            Lấy chi tiết bản nộp kèm thông tin siêu dữ liệu (metadata) của Chapter và Task liên quan.
        \item \texttt{findAllChapterSubmissions(): array}\\
            Lấy toàn bộ danh sách nộp bản thảo Chapter trong hệ thống.
    \end{itemize}

    \item \textbf{Lớp Review:}
    \begin{itemize}
        \item \texttt{findBySubmissionId(submissionId: int): array}\\
            Lấy danh sách nhận xét cho một lần nộp bài cụ thể.
        \item \texttt{countReviewsByEditor(editorId: int, status: string): int}\\
            Đếm số lượng đánh giá của Editor theo trạng thái bài nộp.
        \item \texttt{findByReviewerId(reviewerId: int): array}\\
            Lấy danh sách các phản hồi, đánh giá được thực hiện bởi người dùng.
    \end{itemize}

    \item \textbf{Lớp EditorAnnotation:}
    \begin{itemize}
        \item \texttt{findByPageId(pageId: int): array}\\
            Lấy danh sách các ghi chú báo lỗi của Editor trên trang truyện kèm theo tên Editor.
        \item \texttt{deleteByPageId(pageId: int): bool}\\
            Xóa toàn bộ các ghi chú báo lỗi trên một trang truyện cụ thể.
    \end{itemize}

    \item \textbf{Lớp SubmissionAnnotation:}
    \begin{itemize}
        \item \texttt{findBySubmissionId(submissionId: int): array}\\
            Lấy danh sách ghi chú lỗi trực quan trên bản nộp của trợ lý kèm tên đầy đủ người đánh giá.
    \end{itemize}

    \item \textbf{Lớp BoardVote:}
    \begin{itemize}
        \item \texttt{castVote(seriesId: int, memberId: int, vote: string): bool}\\
            Ghi nhận lượt bỏ phiếu của thành viên hội đồng cho đề xuất truyện.
        \item \texttt{getMemberVote(seriesId: int, memberId: int): string}\\
            Lấy ý kiến bỏ phiếu hiện tại của một thành viên cụ thể.
        \item \texttt{getApprovalStats(seriesId: int): array}\\
            Tính toán động tỉ lệ phần trăm tán thành và số lượng phiếu bầu đạt được.
    \end{itemize}

    \item \textbf{Lớp SeriesRanking:}
    \begin{itemize}
        \item \texttt{findBySeriesId(seriesId: int): array}\\
            Lấy danh sách lịch sử điểm số và thứ hạng của bộ truyện.
        \item \texttt{getSeriesCountByPeriod(period: string): int}\\
            Đếm số lượng bộ truyện đã xếp hạng trong một kỳ đánh giá cụ thể.
        \item \texttt{getTopSeriesByPeriod(period: string, limit: int): array}\\
            Lấy danh sách các bộ truyện có điểm cao nhất trong kỳ đánh giá.
        \item \texttt{getBottomSeriesByPeriod(period: string, limit: int): array}\\
            Lấy danh sách các bộ truyện có điểm thấp nhất trong kỳ đánh giá.
        \item \texttt{getLatestPeriod(): string}\\
            Lấy kỳ xếp hạng gần nhất hiện có trên hệ thống.
        \item \texttt{getPreviousRanking(seriesId: int, period: string): array}\\
            Lấy thông tin thứ hạng và điểm số của bộ truyện trong kỳ đánh giá trước đó.
        \item \texttt{checkDuplicateRanking(seriesId: int, period: string): bool}\\
            Kiểm tra xem bộ truyện đã có bản ghi xếp hạng trong kỳ tương ứng chưa.
        \item \texttt{findByMangakaId(mangakaId: int): array}\\
            Lấy lịch sử xếp hạng của các bộ truyện thuộc Mangaka sở hữu.
        \item \texttt{getLatestRanking(): array}\\
            Lấy danh sách xếp hạng mới nhất của các bộ truyện.
        \item \texttt{findWithDetails(id: int): array}\\
            Lấy chi tiết bản xếp hạng kèm thông tin bộ truyện và ban biên tập liên quan.
        \item \texttt{findAllWithSeries(): array}\\
            Lấy toàn bộ lịch sử xếp hạng tất cả các bộ truyện trên hệ thống.
    \end{itemize}

    \item \textbf{Lớp Notification:}
    \begin{itemize}
        \item \texttt{createNotification(userId: int, type: string, message: string, relatedId: int = null): bool}\\
            Tạo và gửi thông báo mới cho người dùng.
        \item \texttt{findByUserId(userId: int): array}\\
            Lấy toàn bộ danh sách thông báo của người dùng.
        \item \texttt{findUnreadByUserId(userId: int): array}\\
            Lấy danh sách các thông báo chưa đọc của người dùng.
        \item \texttt{getUnreadCount(userId: int): int}\\
            Đếm số lượng thông báo chưa đọc của người dùng.
        \item \texttt{getLatestNotifications(userId: int, limit: int = 5): array}\\
            Lấy danh sách thông báo mới nhất phục vụ hiển thị nhanh (dropdown).
        \item \texttt{markAsRead(notificationId: int, userId: int): bool}\\
            Đánh dấu một thông báo cụ thể là đã đọc.
        \item \texttt{markAllAsRead(userId: int): bool}\\
            Đánh dấu toàn bộ thông báo của người dùng là đã đọc.
        \item \texttt{findAllWithUser(): array}\\
            Lấy danh sách toàn bộ thông báo của mọi người dùng kèm tên chi tiết.
        \item \texttt{getTypeDetails(type: string): array}\\
            Lấy thông tin chi tiết (badge màu, tên hiển thị) theo từng loại thông báo.
    \end{itemize}

    \item \textbf{Lớp SystemLog:}
    \begin{itemize}
        \item \texttt{logAction(userId: int, action: string, details: string): bool}\\
            Ghi nhận nhật ký hoạt động nghiệp vụ của người dùng vào hệ thống.
        \item \texttt{getPaginatedLogs(limit: int, offset: int): array}\\
            Lấy danh sách log có phân trang kèm theo thông tin chi tiết người dùng thực hiện.
    \end{itemize}
\end{itemize}

\subsection{Thiết kế biểu đồ trạng thái (State Machine Diagram)}
\label{subsec:sdd_state_machine}

Biểu đồ trạng thái (State Machine Diagram) được sử dụng để mô tả chi tiết vòng đời hoạt động và các sự kiện chuyển trạng thái của các thực thể cốt lõi trong hệ thống:

\subsubsection{Biểu đồ trạng thái Chương truyện (Chapter)}
\label{subsubsec:state_chapter}
Mô tả vòng đời của chương truyện từ khi là bản thảo thô cho đến khi xuất bản chính thức.
\begin{figure}[H]
    \centering
    \includegraphics[width=0.9\textwidth]{Chapter3_IMAGE/State_Machine_Chapter.png}
    \caption{Biểu đồ trạng thái của Chương truyện (Chapter)}
    \label{fig:state_chapter}
\end{figure}

\subsubsection{Biểu đồ trạng thái Trang truyện (Page)}
\label{subsubsec:state_page}
Mô tả sự dịch chuyển trạng thái vẽ của một trang Manga cụ thể khi có sự tác động phân công công việc.
\begin{figure}[H]
    \centering
    \includegraphics[width=0.8\textwidth]{Chapter3_IMAGE/State_Machine_Page.png}
    \caption{Biểu đồ trạng thái của Trang truyện (Page)}
    \label{fig:state_page}
\end{figure}

\subsubsection{Biểu đồ trạng thái Phân vùng trang truyện (PageRegion)}
\label{subsubsec:state_pageregion}
Mô tả vòng đời của một phân vùng ảnh vẽ tay nhỏ trên trang truyện từ khi được đánh dấu cho đến khi hoàn thành.
\begin{figure}[H]
    \centering
    \includegraphics[width=0.8\textwidth]{Chapter3_IMAGE/State_Machine_Page_Region.png}
    \caption{Biểu đồ trạng thái của Phân vùng trang truyện (PageRegion)}
    \label{fig:state_pageregion}
\end{figure}

\subsubsection{Biểu đồ trạng thái Bộ truyện Manga (Series)}
\label{subsubsec:state_series}
Mô tả vòng đời của dự án truyện Manga từ đề xuất kế hoạch, sáng tác thực tế đến khi kết thúc hoặc đình bản.
\begin{figure}[H]
    \centering
    \includegraphics[width=0.8\textwidth]{Chapter3_IMAGE/State_Machine_Series.png}
    \caption{Biểu đồ trạng thái của Bộ truyện (Series)}
    \label{fig:state_series}
\end{figure}

\subsubsection{Biểu đồ trạng thái Công việc giao trợ lý (Task)}
\label{subsubsec:state_task}
Mô tả tiến độ phân công, tiếp nhận và hoàn thiện các Task vẽ hỗ trợ của các Assistant.
\begin{figure}[H]
    \centering
    \includegraphics[width=0.8\textwidth]{Chapter3_IMAGE/State_Machine_Task.png}
    \caption{Biểu đồ trạng thái của Công việc (Task)}
    \label{fig:state_task}
\end{figure}

\subsubsection{Class Diagram tổng thể}
\label{subsubsec:class_overview}

Class Diagram tổng thể thể hiện toàn bộ các lớp cốt lõi của hệ thống cùng các mối quan hệ liên kết, kế thừa và phụ thuộc giữa chúng.

\begin{figure}[H]
    \centering
    \includegraphics[width=\textwidth, height=0.9\textheight, keepaspectratio]{Chapter3_IMAGE/class_diagram.png}
    \caption{Class Diagram tổng thể}
    \label{fig:class_diagram_tong_the}
\end{figure}

\subsubsection{Mô tả các mối quan hệ}
\label{subsubsec:class_relationships}

Các mối quan hệ (Association/Multiplicity) giữa các lớp trong Class Diagram thể hiện luồng nghiệp vụ và tương tác dữ liệu cốt lõi của hệ thống, bao gồm:

\begin{itemize}
    \item \textbf{Role (1) - (*) User:} Một vai trò (VD: Admin, Manager, Assistant, Editor) có thể được gán cho nhiều người dùng. Mỗi người dùng bắt buộc thuộc về một vai trò nhất định để phân quyền hệ thống.
    \item \textbf{User (1) - (*) Series:} Một tác giả (Mangaka) có thể sáng tác và quản lý nhiều bộ truyện (Series) khác nhau.
    \item \textbf{Series (1) - (*) Chapter:} Một bộ truyện bao gồm nhiều chương (Chapter) nối tiếp nhau.
    \item \textbf{Chapter (1) - (*) Page:} Mỗi chương truyện bao gồm nhiều trang (Page) cụ thể chứa hình ảnh nội dung truyện.
    \item \textbf{User (1) - (*) Task:} Người dùng đóng vai trò giao việc (Manager/Mangaka) có thể phân công nhiều công việc (Task). Đồng thời, người dùng đóng vai trò nhận việc (Assistant) có thể tiếp nhận và thực hiện nhiều Task.
    \item \textbf{Task (1) - (*) Submission:} Một công việc được giao có thể tạo ra nhiều bản nộp (Submission) khác nhau nếu yêu cầu phải nộp lại nhiều lần.
    \item \textbf{Submission (1) - (*) Review:} Mỗi bản nộp của Assistant có thể nhận được nhiều lượt đánh giá, phản hồi (Review) từ cấp trên để chỉnh sửa cho hoàn thiện.
    \item \textbf{Submission (1) - (*) SubmissionAnnotation:} Mỗi bản nộp (Submission) của Trợ lý có thể nhận được nhiều ghi chú lỗi trực quan (SubmissionAnnotation) từ Mangaka để phục vụ chỉnh sửa.
    \item \textbf{User (1) - (*) SubmissionAnnotation:} Tác giả (Mangaka) có thể tạo nhiều ghi chú lỗi trực quan trên các bản nộp khác nhau.
    \item \textbf{Series (1) - (*) SeriesRanking:} Mỗi bộ truyện sẽ được đánh giá và lưu lịch sử qua nhiều kỳ xếp hạng (SeriesRanking) khác nhau dựa trên điểm số (score) và phiếu bình chọn (votes) thu thập được.
    \item \textbf{User (1) - (*) SeriesRanking:} Ban biên tập (User) là những người thực hiện đánh giá và cập nhật xếp hạng cho các bộ truyện.
    \item \textbf{User (1) - (*) Notification:} Mỗi người dùng sẽ có hộp thư chứa nhiều thông báo (Notification) từ hệ thống để cập nhật tiến độ công việc hoặc kết quả duyệt bài.
\end{itemize}

